\documentclass[../../main.tex]{subfiles}
\begin{document}

\begin{center}
  {\fontsize{48pt}{58pt}\selectfont \textbf{1978}\normalsize\textbf{年度}}
\end{center}

\begin{table}[htb]
  \centering
  \caption{出題テーマと難易度}
  \begin{tabular}{|c|p{9cm}|c|}
    \hline
    \textbf{問題番号} & \textbf{テーマ} & \textbf{難易度} \\
    \hline
    第1問 & 円と平行四辺形で囲まれる領域の面積 & \\
    \hline
    第2問 & 区間上の四次関数の最大値とグラフ & \\
    \hline
    第3問 & 放物線の法線の本数と領域の面積 & \\
    \hline
    第4問 & 行列の累乗とベクトルの極限方向 & \\
    \hline
    第5問 & 三角形内の距離の二乗和と区分求積 & \\
    \hline
    第6問 & 放物線内の最大面積三角形 & \\
    \hline
  \end{tabular}
\end{table}

\begin{table}[htb]
  \centering
  \caption{解答}
  \renewcommand{\arraystretch}{1.4}
  \begin{tabular}{|c|c|p{9cm}|}
    \hline
    \textbf{問題番号} & \textbf{小問} & \textbf{答え} \\
    \hline
    第1問 & -- & $OPCQ$は平行四辺形．面積は$5\pi-21+\left(12-\dfrac{8\pi}{3}\right)\sqrt{3}$ \\
    \hline
    第2問 & -- & $-3\le t\le-1$では$g(t)=f(t+1)$，$-1\le t\le0$では$36$，$0\le t\le2$では$f(t)$，$2\le t\le3$では$f(t+1)$．最大値は$84$（グラフは解答参照） \\
    \hline
    第3問 & (1) & 点$P$の領域は解答の図を参照 \\
    \cline{2-3}
    & (2) & 面積は$\dfrac{88\sqrt{2}}{135}$ \\
    \hline
    第4問 & -- & $A^n=
    \begin{pmatrix}(1/3)^n & \dfrac{15}{8}(3^n-(1/3)^n)\\0&3^n
    \end{pmatrix}$．$(u,v)$は$b=0$なら$\left(\dfrac{a}{|a|},0\right)$，$b\ne0$なら$\left(\dfrac{15b}{17|b|},\dfrac{8b}{17|b|}\right)$ \\
    \hline
    第5問 & -- & 極限値は$\dfrac{-a^2+3b^2+3c^2}{6}$ \\
    \hline
    第6問 & (1) & $0<x\le\dfrac{2}{3}$ \\
    \cline{2-3}
    & (2) & $x\le\dfrac{2}{3}$なら$f(x)=\dfrac{(2x-x^2)(1-x)}{2}$，$x\ge\dfrac{2}{3}$なら$f(x)=\dfrac{(2-x)x^2}{4}$．極大値$\dfrac{\sqrt{3}}{9}$，極小値$\dfrac{4}{27}$（グラフは解答参照） \\
    \hline
  \end{tabular}
\end{table}

\paragraph{試験形態}
6問．試験時間は要確認．

\end{document}
